\documentclass{article}
\usepackage{spconf,amsmath,graphicx}

\usepackage[T1]{fontenc}
\usepackage{lmodern}
\usepackage{amssymb,booktabs,array,tabularx}
\usepackage{enumitem}
\usepackage[hidelinks]{hyperref}
\usepackage{url}
\graphicspath{{../results/figures/}}

\renewcommand{\topfraction}{0.92}
\renewcommand{\bottomfraction}{0.85}
\renewcommand{\textfraction}{0.07}
\renewcommand{\floatpagefraction}{0.75}
\setcounter{topnumber}{2}
\setcounter{totalnumber}{3}
\setlist[enumerate]{nosep,leftmargin=1.2em,labelsep=0.35em,topsep=2pt}
\setlist[itemize]{nosep,leftmargin=1.0em,labelsep=0.35em,topsep=2pt}

\title{Does a Music Transformer Know What Key It Is In?\\
Causal Interventions on Emergent Tonal Representations}

\name{Genki Masuyama$^{\dagger}$, Keigo Sakurai$^{\dagger\dagger}$,
Ren Togo$^{\dagger\dagger}$, Takahiro Ogawa$^{\dagger\dagger}$,
Miki Haseyama$^{\dagger\dagger}$}

\address{
$^{\dagger}$Graduate School of Information Science and Technology, Hokkaido University, Japan\\
$^{\dagger\dagger}$Faculty of Information Science and Technology, Hokkaido University, Japan\\
E-mail: \{masuyama, sakurai, togo, ogawa, mhaseyama\}@lmd.ist.hokudai.ac.jp
}

\begin{document}
\ninept
\maketitle
\begin{abstract}
% 小さな分類器は音楽モデルの活性から調を読み出せるが，それは，生成の間にモデルがその調を使っていることを示さない．
% 我々はモデルに調のラベルを与えず，語彙も調の記号を持たない記号音楽で，Transformer を訓練する．
% 生成の途中で，1つの層において，調を担う部分活性を，別の調の部分活性に置き換える．入力は変えないので，出力に生じた変化はこの置き換えだけによる．
% するとその後の出力のおよそ3本に1本が，書き込んだ調にちょうど変化する；ランダムな方向へ書き込んだ場合は 25 本に1本である．短調では効果はさらに大きい．
% モデルに音を生成させる前の時点で，これから選ぼうとしている音がすでに書き込んだ調のものになっている．つまりモデルは，その調を自分の出力から読み取っているのではない．
% どちらの結果も，実音楽で訓練された公開モデルにそのまま当てはまる．
% したがって音楽 Transformer は調を計算し，我々が書き換える形で保持し，生成しながらそれを参照している．
A small classifier can read the musical key from a music model's activations, but that does not show that the model uses the key as it generates.
We train Transformers on symbolic music, giving the model no key labels and a vocabulary that contains no key symbols.
During generation we replace, at one layer, the key-carrying part of the activation with the same part from a different key, and we leave the input unchanged, so only that replacement can change the output.
About one continuation in three then moves into exactly the key we install, against one in twenty-five when we install along random directions instead; with minor keys the effect is larger still.
Before we let the model generate a note, the notes it is about to choose already belong to the installed key, so the model is not reading that key from its own output.
Both results carry over to a public model trained on real music.
A music Transformer therefore computes its key, holds it in a form we can replace, and consults it as it generates.
\end{abstract}

\begin{keywords}
music generation, interpretability, activation patching, causal intervention, tonality, probing
\end{keywords}

% I. はじめに
\section{Introduction}
% 音楽生成モデルは，いま自分がどの調にいるのかを知っているかのように振る舞う．
% 調とは，ある一節が「ホーム」として扱う少数の音高の集まりである．音楽はそこに落ち着き，そこを離れ，また戻ってくる．
% 学習済みモデルが生成する音楽でも同じことが起きる．途中で別の調へ移り，また元の調へ戻ってくることが多い．
% モデルはこれをあまりに自然にこなすため，そこでひとつの疑問が生まれる．モデルは内部のどこかで調を追い続けているのか，それとも音符の表面的なパターンをなぞっているだけなのか．
Music generation models behave as if they knew what key they were in~\cite{hadjeres2017deepbach,huang2019musictransformer,copet2023musicgen}.
A key is the small set of pitches a passage treats as home: the music settles on them, leaves them, and comes back---and generated music does the same, moving away partway and returning.
It happens so naturally that it raises a question: does the model track the key somewhere inside itself, or only reproduce surface patterns of notes?


% 手元にある証拠はこの問いに答えられない．どの証拠も，我々が設定した調にモデルを従わせるものではないからである．
% 既存研究がモデルの内部に働きかける仕方は2つある．
% 1つ目は読むことである：プロービング研究は活性の上で小さな分類器を訓練し，音符や和音から音階までの音楽理論の概念を復元し，曲の調も読み出してきた．
% しかしプローブが示すのは調が「読める」ことまでであり，プローブは，モデルがその信号を使っていなくても成功しうる．
% 2つ目は押すことである：ステアリング研究は，言語モデルにならって学習したベクトルを活性に足し，出力の属性を動かす．
% しかし足したベクトルは，活性のうち現在の調を担う部分をそのまま残すので，ホ長調の方向へ押してもモデルの調がホ長調に設定されることはない．
% 要するに，調を読むことも出力を押すこともできるが，調を設定することはまだできない．
The evidence we have cannot answer that, because none of it makes the model follow a key that we set.
Existing work reaches inside a music model in two ways.
Probing trains a small classifier on the activations and recovers music-theory concepts~\cite{syntheory2024}, the key among them~\cite{castellon2021calm} --- but a probe shows only that the key is readable, and can succeed on a signal the model ignores~\cite{hewitt2019control,belinkov2022probing}.
Steering adds a learned vector to the activations, following language models~\cite{li2023iti,zou2023repe}, and an output attribute moves~\cite{smitin2024,facchiano2025patching,singh2026discovering,musicrfm2026} --- but the added vector leaves the part of the activation carrying the current key in place, so pushing toward E major competes with the key that is still written there; what that preserved component costs is something we measure (Table~\ref{tab:main}, bottom).
% 「設定できない」は held-out の C 条件（加算・2倍変位で 0.471 > install 0.355）と矛盾するので使えない．
% 実際の空白は「同じ摂動予算のもとで，旧調成分を消すことに独立の価値があるか」が測られていないこと．
We can read the key and push the output, but no one has asked what it costs to \emph{replace} what is written there rather than add to it.

% 関連研究セクションがないので，ここで研究をする意味をさらに補強するために，先行研究を入れ，問いを明確化
% 活性を上書きすればこの種の問いは決着するが，音楽モデルで調を上書きした研究はまだない．
% ゲームも同じ隔たりに直面し，それを埋めたのは上書きだった：オセロでは，手の予測だけで訓練されたモデルが盤面の内部表現を作り，研究者がその表現を偽の盤面で上書きし，入力には手を触れずにおくと，モデルは偽の盤面のほうで合法な手を指した．
% この種の編集は activation patching と呼ばれ，その後ほかの領域でも標準的なテストになっている．
% 音楽では，最も近い2つの試みも，どちらも調を書き込んではいない：Ma と Xia は和音表現を上書きしており，Vélez Vásquez らは MusicGen から概念を除去したが，別の概念を入れていない．
% 音楽モデルに調を書き込んだ研究はまだないので，我々の問いの操作的な形は未解決である．
Overwriting an activation settles questions of this kind, and it closed the same gap in games~\cite{mcgrath2022chess}: an Othello model trained only on moves builds a board representation~\cite{li2023othello,nanda2023emergent}, and overwriting it with a fake board --- input untouched --- makes the model play the fake board's legal moves.
Such edits, activation patching~\cite{geiger2021causal,meng2022rome}, are now a standard test elsewhere~\cite{karvonen2024chess,maze2024,iris2026}.
In music the two nearest attempts write no key: Ma and Xia~\cite{ma2024root} overwrite chords --- a bar-scale attribute, where a key governs whole passages --- and V{\'e}lez V{\'a}squez et al.~\cite{velez2024exploring} remove a concept from MusicGen without adding another; the operational form of our question is open.

% 活性パッチングによる介入の概要.活性のうち調を担う成分を目標とする調の成分に置き換え,その他の活性成分と入力は変更しない.
\begin{figure}[!t]
\centering
\includegraphics[width=0.88\columnwidth]{fig1.pdf}
\caption{Overview of the activation-patching intervention.
We replace the key-carrying component of an activation with that of a target key while leaving the remaining activation and the input unchanged.}
\label{fig:overview}
\end{figure}

% 我々はこの問いを，モデルの内部で調を変え，その後に何を書くかを見ることで決着させる（図1）．
% テストの核にあるのは1つの操作であり，それをモデルが生成する各位置に適用する：1つの層において，調を担う部分活性を，別の調の部分活性に置き換え，入力は変えない．だから出力に生じた変化はこの置き換えだけによる．
% これは activation patching を部分空間に限定したもの——distributed interchange intervention——であり，以後この操作を「調を install する（書き込む）」と呼ぶ．
% 音楽には，編集が効いたかをすぐに判定する仕組みがない．そこで効いてくるのが調である：どの調にいるかは確立された基準が決めており，語彙は調の記号を持たず，編集の成否は音符を数えれば確かめられる．
% 書き込んだ値に音楽が従えば，モデルはその値を参照していたことになる．
% 従わなかった場合に分かるのは，モデルがその層のその部分活性を使っていない，ということまでであり，調を持たないと結論することはできない．
We settle the question by changing the key inside the model and watching what it writes next (Fig.~\ref{fig:overview}): at every position the model generates, we replace the key-carrying part of one layer's activation with the same part from a different key, leaving the input unchanged, so only the replacement can change the output --- activation patching restricted to a subspace, a distributed interchange intervention~\cite{geiger2024das}; we say we \emph{install} the key.
Key is the right concept to test because an established criterion decides which key a passage is in~\cite{krumhansl1982,temperley1999ks}, our vocabulary holds no key symbols, and counting notes checks the edit.
If the music follows the installed value, the model was consulting it; if not, we learn only that it does not use that part of that layer's activation --- not that it has no key.

% 本論文では，新しい生成手法ではなく分析を提示する：音楽を生成する新しい方法は提案せず，生成品質の向上も主張しない．
% その代わりに我々が提供するのは，誰でも再実行できるテストである：データ，コード，そして事前に固定した手続きを公開する．
% 本論文の貢献は次の2点である．
% ・我々は音楽 Transformer の活性をプローブし，モデルが調を入力の音符から写しているのではなく，自ら計算していることを示す．
% ・我々は生成の途中でモデル内部に別の調を書き込み，音楽がそれに従うことを，自作モデルでも，実音楽で訓練された公開モデルでも示す．
This is an analysis, not a generation method: we claim no quality gain, and we release the data, code, and pre-registered protocol so the test can be rerun.
Our contributions are twofold:
\begin{itemize}
\item We probe the activations of a music Transformer and show that it computes the key itself, rather than copying it from the notes in its input.
\item We install a different key inside the model during generation and show that the music follows it, both in our own models and in a public model trained on real music.
\end{itemize}

% II. 我々のテスト
\section{Probing and Causal Intervention}
\label{sec:method}
% 本研究は問いを2段階に分けて尋ねる．
% 実験1は，モデルがそもそも調を計算しているのかを尋ねる：調は，音符からよりも活性からのほうが読み取りやすいのか．
% 実験2は，モデルが書くときにその調を参照しているのかを尋ねる：それを置き換えたら，音楽はそれに従うのか．
% 本節では2つの実験を述べる；データとモデル，測るもの，そして因果の主張を確かなものに保つための「探索してから検証する」手続きは III 節で述べる．
% 記法として，層 ℓ・位置 t における活性を h_ℓ(t)，そこでの真の調を κ(t) と書く．κ(t) は 24 クラス（12 の主音それぞれに長調と短調）を取る．
% 層 0–7 は8つの Transformer ブロックの出力を指す；トークン埋め込みはまだ何の計算も経ていないので，層として数えない．
% プローブは 24 クラスすべてを判別する；編集は，主テストでは 12 の長調を，副テストでは 12 の短調を書き込む．
Experiment~1 asks whether the model computes a key at all: is the key easier to read from the activations than from the notes?
Experiment~2 asks whether it consults that key when it writes: if we replace it, does the music follow?
Section~\ref{sec:setup} gives the data, models, measures, and the search-then-test protocol that keeps the causal claim honest.
Write $h_\ell(t)$ for the activation at layer $\ell$ and position $t$ (layers index the Transformer block outputs), and $\kappa(t)$ for the true key there, one of $24$ classes ($12$ tonics $\times$ major/minor).
The probe distinguishes all $24$; the edits install the $12$ major keys in the primary test and the $12$ minor keys in a secondary one.

\subsection{Experiment 1: Reading the Key with a Probe}
\label{sec:exp1}
% 実験1は，調が音符そのものからよりも h_ℓ(t) からのほうが読み取りやすいかを尋ねる．
% 活性は音符から計算されるので，調の情報が中にあること自体は当然である．問うているのは，それが線形に読み出せる形になっているかである．
% h_ℓ(t) から κ(t) を予測する線形プローブを訓練し，訓練とテストは曲単位で分ける．
% プローブの得点だけでは何も言えないので，音符から調を当てる比較対象を置く：各位置で直前のトークンをさかのぼり，窓の中の音高トークンについて 12 のピッチクラス（C，C♯，…，B）がそれぞれ何回現れたかを数え，その 12 個の数からその位置の調を当てる——回数へのロジスティック回帰と，ピッチクラスの出現回数から調を推定する標準手法である Krumhansl–Schmuckler（KS）法の2通りで．
% 窓を広げれば音符数えが勝てるのではないかという疑いに答えるため，窓の取り方を変えた追加実験も行い，そのうち最も強いものを音符数えベースラインと呼ぶ（窓の設定は補足資料）．
% プローブがこれを上回れば，モデルは音符が与えるものに何かを足したことになる．
% さらに，モデルが調について何も計算していなくても得点が出る経路が2つあるので，その分を測って差し引く．
% ・プローブの強さ：強力なプローブは意味のないラベルからでも点を取る．そこで control task を走らせる：各曲の調をランダムな別の調に置き換えて得点を測り，それを，報告するすべての優位から引く．こうして各優位は selectivity になる．
% ・未訓練のネットワーク：構造だけで得点が出るかもしれないので，同じプローブを未訓練のモデルにも走らせる．
The activation is computed from the notes, so the key is in there in some form by default; the question is whether a linear probe can read it.
We train one~\cite{alain2017probes,belinkov2022probing} to predict $\kappa(t)$ from $h_\ell(t)$, splitting train and test by piece.
A probe score alone says nothing, so we compare it against a baseline that reads the key from the notes: at each position we look back over the preceding tokens, count how many times each of the twelve pitch classes (C, C$\sharp$, \dots, B) occurs among the pitch tokens in the window, and predict the key from those twelve counts, by logistic regression and by Krumhansl--Schmuckler (KS) matching~\cite{krumhansl1990,temperley1999ks}, the standard method for estimating key from pitch-class counts.
To answer the objection that a note counter might win with more history, we also ran the counter with wider windows, and we call the strongest of all these baselines the note-counting baseline; the supplementary gives the window settings.
If the probe beats it, the model has added something to what the notes already give.
Two further routes could produce a probe score even when nothing about the key is being computed, so we measure what each contributes and subtract it.
\begin{itemize}
\item \textbf{Probe power.} A powerful probe scores points even on meaningless labels. We therefore run a control task~\cite{hewitt2019control}: we replace the key of each piece with a random one, measure the probe there, and subtract that score from every margin we report, so each margin is a selectivity.
\item \textbf{An untrained network.} Its activations might already score something from structure alone, so we run the same probe on one.
\end{itemize}

\subsection{Experiment 2: Installing a Key During Generation}
\label{sec:exp2}
% 実験2（本論文の主実験）は，生成中に別の調を書き込んだとき音楽がそれに従うかを尋ねる．
% 1回の実行は5つの手順で進む．
Experiment~2, the main one, asks whether the music follows the key when we install a different one during generation; one run proceeds in five steps.
\begin{enumerate}
% 1．部分空間を選ぶ：調は活性の全体を占めているわけではないので，まず調を担う部分を特定する．方向の作り方は3通り試した：プローブの重み，調ごとの平均活性，そして編集の効果そのものを目的関数にして学習した方向（distributed alignment search）である．3通り試すのは，編集が失敗したとき，モデルが調を参照していないのか，方向の見つけ方が悪かったのかを切り分けるためである．いずれの場合も，得られた 24 本——調のクラス1つにつき1本——を行列 V に並べ，特異値分解で正規直交化する．こうすると P_V=VV^T は，活性を調の部分だけに射影し，残りには触れない．探索段階で3通りを比べ，プローブの重みが最も良かった（III-D 節）．
\item \textbf{Choose the subspace.} The key does not occupy the whole activation, so we first locate the part that carries it, trying three kinds of directions so a failure could not be blamed on one poor choice: the probe weights, the per-key mean activations, and directions trained on the edit's own effect (distributed alignment search~\cite{geiger2024das}). The $24$ directions---one per key class---are orthonormalized as the columns of $V$, so $P_V{=}VV^{\!\top}$ projects onto that part alone. The search stage compared the three; the probe weights won (Sec.~\ref{sec:confirm}).
% 2．書き込む値を選ぶ：真の調が標的 κ* である位置すべてにわたって層 ℓ の活性を平均し，μ_{κ*} を得る．これはその調にいるときの活性の典型である．
\item \textbf{Choose the value to install:} $\mu_{\kappa^{*}}$, the layer-$\ell$ activation averaged over every position whose true key is the target --- a typical moment in that key.
% 3．プロンプト：調が1つに定まった8小節を与える．
\item \textbf{Prompt} the model with $8$ bars in a single key.
% 4．書き込む：モデルが生成を進めるあいだ，第9小節以降の各位置について，層 ℓ の活性が次のブロックへ渡る前に式(1)を適用する．他の層と，プロンプトの位置には手を加えない．
%    最初の2項が活性にある調を消し，第3項が標的の調をその場所に書き込むので，部分空間の外側は何も変わらない．
\item \textbf{Install:} apply Eq.~\eqref{eq:edit} to the layer-$\ell$ activation at each generated position from bar~$9$ on, before it passes to the next block; no other layer, no prompt position.
\begin{equation}
h \;\leftarrow\; h - P_V h + P_V\,\mu_{\kappa^{*}}.
\label{eq:edit}
\end{equation}
The first two terms delete the key that is in the activation and the third writes the target key in its place, so nothing outside the subspace changes.
% 5．サンプリングして採点する：プロンプト1本・標的1調につき継続を1本，温度 1.0・top-p 0.95・固定シードで最大 384 トークン生成し，第9–24小節を採点する．
\item \textbf{Sample and score} one continuation per prompt and target key---temperature $1.0$, top-$p$ $0.95$, a fixed seed, at most $384$ tokens---and score bars $9$--$24$.
\end{enumerate}

% 生成を伴わない読み取り．
% 上の手順5は，編集の効果を16小節の音楽で測る．同じ編集は，最初の1手だけで測る変種にも使える．
% 調に中立な小節記号と拍位置を我々自身が付け足し，第9小節の境界で編集を適用し，音を1つも生成させずに次の音の確率をそのまま読む（IV-C 節）．
% これにより，編集直後の最初の決定を，モデルが本来書いていたはずの音符の影響から切り離せる．
\noindent\textbf{A reading without generation.}
The same edit supports a variant that measures only the first decision: we append a key-neutral bar mark and beat position ourselves and read the next-note probabilities with nothing generated (Sec.~\ref{sec:nextpitch}), isolating that decision from any note the model wrote.

% III. 実験の構成 
\section{Experimental Setup}
\label{sec:setup}

\subsection{Data}
% このテストには，全音符の真の調が必要である．
% 実音楽でこれに最も近いのは機能和声の分析コーパスだが，そのラベルは人間の分析であり，規模も数百曲にとどまるので，言語モデルの訓練には足りない．
% そこでデータは自分たちで生成する——SynTheory のような診断用データセットの設計に倣う．
% 音楽の教科書にある規則に従って，プログラムがまず調を1つ選び，その調に合う和音を順に並べ，各和音を4つの音として書き出し，半数の曲では旋律を上に加える．
% 訓練用に 20万曲，検証用とテスト用に各1万曲を生成し，各曲は 278–508 トークンで 0–3 回の調の変化を含む．
% プログラムは調を決めてから音符を書くので，全トークンの調が確定する：ラベルは推定ではなく記録である．
% さらに調はほぼ均等に現れる：テスト分割では，最も多い調のトークン数は最も少ない調の 1.12 倍にとどまる．
% モデルはこのラベルを決して見ない；使うのはモデルの外側だけである：プローブの訓練と採点，そして調ごとに活性を平均して書き込む値を作ることに使う（II-B 節）．
% 語彙は 124 種類のトークン——小節，位置，絶対 MIDI 音高，音の長さ，構造記号——であり，調や和音の記号は含まれない．
% モデルが調を持つに至ったのなら，それはモデルが自分で作ったものである．
Our test needs the true key of every note, and human-analysed corpora such as When in Rome~\cite{gotham2023wir} cover only hundreds of pieces --- too few to train on.
We therefore generate the data ourselves, following the design of diagnostic datasets such as SynTheory~\cite{syntheory2024}.
Following textbook harmony, a program picks a key, writes a chord sequence in it (four notes each, half the pieces with a melody on top): $200{,}000$ training pieces and $10{,}000$ each for validation and test, $278$--$508$ tokens long, $0$--$3$ key changes, keys near-uniform (most/least frequent ratio $1.12$).
The program decides the key before writing the notes, so every token's key is a record, not an estimate; the model never sees these labels, which serve only to train and score the probe and to average each key's activations into the value we write in (Sec.~\ref{sec:exp2}).
The vocabulary holds $124$ token types---bars, positions, absolute MIDI pitches, note lengths, structural marks---and no key or chord symbol.
If the model ends up with a key, it built one.

\subsection{Models}
% 報告する数値はすべて1つのモデルから得たものであり，計6モデルを訓練するのは結果が偶然でないことを確かめるためである．
% どのモデルも 8 層・活性 512 次元・約 25M パラメータの GPT-2 型デコーダで，次トークン予測だけを目標に訓練する．
% 6モデルは2つの軸の組み合わせである．
% 1つ目は移調による増強で，訓練データの各曲を12の調すべてに複製する：調の表現が訓練データの調の偏りに由来するなら，増強を外せば結果が変わるはずである．
% 2つ目は乱数シードで，3種を使う：ある発見が特定の初期値に依存するなら，シードを変えれば結果が変わるはずである．
% 主実験に使うのは，移調増強あり・シード0の1モデルである；残りはプローブの結果の確認に使う（IV-A 節）．
% 6モデルはいずれも，検証データの次トークン top-1 精度が 0.878–0.881 であり，最頻トークンを当て続けるだけの 0.434 を大きく上回る．
% 編集が音楽をどれだけ乱すかは，別に訓練した基準モデルのパープレキシティで測る．
% 基準モデルは同じアーキテクチャだが，別のシードと，他のモデルが見ないデータで訓練する；我々はこのモデルをプローブも編集もしない．
% したがって，どのモデルも自分の出力を採点しない．
We report all our numbers from one model and train six in total to check that the result is not an accident: two augmentation regimes (with and without copying every training piece into all twelve keys --- a key representation inherited from the corpus's key distribution should vanish without it) crossed with three seeds.
Each is a GPT-2 style decoder~\cite{vaswani2017attention,radford2019gpt2} with $8$ layers, $512$-dimensional activations, about $25$M parameters, trained on next-token prediction alone.
The main experiments use the augmented seed-$0$ model; the other five check the probe result (Sec.~\ref{sec:h1}).
All six reach next-token top-$1$ accuracy of $0.878$--$0.881$ on validation, well above the $0.434$ of always guessing the most frequent token.
How far an edit disturbs the music is measured by the perplexity of a separate reference model.
The reference model shares the architecture but is trained on a different seed and on data the other models never see, and we never probe or edit it.
No model therefore grades its own output.

% 合成データで得た結論が実音楽でも成り立つかを確かめるため，我々が訓練していない公開モデル群を調べる（表2）．
% 主となるのは Anticipatory Music Transformer で，原著者が Lakh MIDI，MetaMIDI，FMA で訓練したもの（3サイズ）．
% さらに同じ Lakh MIDI で訓練された，トークン化の異なる2つの公開チェックポイント（compound の MMT と REMI 表現）を加える．
% 評価コーパスは2つ：When in Rome の人手分析による調ラベル付きバッハコラール 300 曲と，
% 人手修正済みの調ラベルを持つポップ 909 曲（POP909-CL）である．
% どのモデルでも，プローブは符号化上「次の音高を選ぶ」位置で読む；音の後で読む場合は補足資料で報告する．
To check that our conclusions hold for real music, we also study public models that we did not train (Table~\ref{tab:wild}).
The main one is the Anticipatory Music Transformer~\cite{thickstun2023anticipatory}, trained by its authors on Lakh MIDI~\cite{raffel2016lakh}, MetaMIDI~\cite{ens2021metamidi}, and FMA~\cite{defferrard2017fma}, in three sizes; we add two further public checkpoints trained on Lakh MIDI with different tokenizations, the compound-token Multitrack Music Transformer (MMT)~\cite{dong2023mmt} and its authors' REMI-representation baseline.
We evaluate on two corpora: $300$ Bach chorales with human key labels from the Roman-numeral analyses of When in Rome~\cite{gotham2023wir}, and $905$ pop songs with human-corrected key labels (POP909-CL~\cite{pop909cl2026}); the supplementary gives editions, splits, a gate that checks the labels against the notes, and one that checks each model's re-implemented tokenizer against the model itself (it detects wrong field structure and a wrong time scale; a uniform offset is a transposition in these vocabularies -- verified, not assumed: the note codes are instrument-major, so an off-by-one moves pitch, and it could reach the instrument field only at pitch $127$, which never occurs in the $1.94$M notes used here -- and is caught instead by vocabulary agreement and decode round-trips).
In every model the probe reads at the position where the encoding says the next pitch is chosen; reading after the note instead is reported in the supplementary.

\subsection{What Counts as Success}
\label{sec:measures}
% 以下はすべて，対応する実験を走らせる前に固定した．
Everything below was fixed before its run.

% 何と比べるか．入力に手を加えないので，出力が変わったならそれは置き換えのせいである．
% しかし置き換えの何が効いたのかは分からない：どの方向に書いたのか，どれだけ遠くへ動かしたのか，どの位置で書いたのか．
% 表1に各比較と，それが消す可能性を示す．
% 中段の行は本文の3語（方向・距離・位置）に1対1対応し，第4の軸（操作：置換か加算か）を追加する．
\noindent\textbf{Controls.}
The input never changes, so any output change comes from the replacement---but not from which part of it: the \emph{directions}, the \emph{distance}, the \emph{positions}, or the \emph{operation} itself (replace against add).
The middle rows of Table~\ref{tab:main} remove these one per row, in that order.

% 主指標．実験1はプローブの優位：control補正後のプローブ得点から音符数えベースラインを引いた差（24クラスの macro-F1）であり，成功は曲単位ブートストラップ 95% CI がゼロを含まないことである．
% 実験2は成功率（SR）である．継続が成功と数えられるのは，次の2つの条件を満たしたときである．
% (i) KS 推定器（II-A 節）が，書き込んだ調とちょうど一致すると判定する．推定器は第9–24小節の音高を見る；音が数個ではピッチクラスの出方が定まらず判定が偶然で決まるので，8音未満の継続は失敗と数える．
% (ii) 音楽が壊れていない．調が変わっても壊れていては意味がなく，またモデルを分布外へ飛ばす編集はその回路について何も語らないからである．上限は本物の転調が決める：検証データの自然な転調 7,989 件について，転調前 24 トークンから後 24 トークンへの基準モデルのパープレキシティの上昇を測り，その 90 パーセンタイルが 0.613 であった．編集した継続は，同じプロンプトの編集なしの継続に対し，これ以上パープレキシティを上げてはならない．
% どちらかの条件を満たさない継続は，分母に残して失敗と数える．当てずっぽうは 1/12 である．
% 公開モデルでは条件 (ii) の上限をコラールで測り直し，基準には別の公開チェックポイントを使う．
\noindent\textbf{Primary measures.}
For Experiment~1, the \emph{probe margin}: the control-corrected probe score minus the note-counting baseline, in macro-$F_1$ over the $24$ classes, with success a piece-level bootstrap $95\%$ CI~\cite{efron1987bca} that excludes zero.
For Experiment~2, the \emph{success rate} (SR).
A continuation counts as a success when it meets two conditions.
% (i) KS は継続全体（最大16小節・EOSで打ち切り）を1つの窓として1回だけ推定する．小節別の推定・集約はしない．
% トークン上限 384 の下で実際に生成されるのは 4–14 小節（最終テスト実測，中央値 10）で，8音未満は一度も起きなかった．
(i)~The KS estimator (Sec.~\ref{sec:exp1}) places it in exactly the installed key.
The estimator runs ONCE over all pitches the model generated, from bar~$9$ up to at most bar~$24$ (the token budget yields $4$--$14$ generated bars in practice); a continuation with fewer than $8$ pitches would count as a failure outright, though none occurred in the final test.
(ii)~The music is not broken---an edit that drives the model off distribution says nothing about its circuitry---and real key changes set the bar: the reference model's perplexity rise across $7{,}989$ natural key changes in validation has a $90$th percentile of $0.613$, and an edited continuation may cost at most that much over the unedited continuation of the same prompt.
A continuation that fails either condition stays in the denominator as a failure.
Chance is one key in $12$.
For the public model we measure the bar of~(ii) on the chorales and use a different public checkpoint as the reference.

% 副指標．音階内音符の割合：継続の音のうち与えられた音階に属するもの（オクターブ不問）である．
% 推定器を使わないので，その混同を引き継がない．
\noindent\textbf{Secondary measure.}
The \emph{in-key note share}: continuation notes belonging to a given scale, any octave---estimator-free, so it inherits no estimator confusion.

% 検定．各標的調について，同じプロンプト上で編集がベースラインに勝つかを尋ねる：プロンプトで対にした片側 Wilcoxon 符号順位検定，12調にわたる Holm 補正，効果量 r は順位双列相関である．
\noindent\textbf{Tests.}
Per target key, a one-sided Wilcoxon signed-rank test paired by prompt (does the edit beat its baseline on the same prompts?), Holm-corrected across the $12$ keys~\cite{holm1979}; $r$ is the rank-biserial correlation.

\subsection{Search Stage and Final Test}
\label{sec:confirm}
% II-B 節の手順には決めるべきことが多く残っている——どの層か，どの方向か，どの位置か——そして報告する数値を出すのと同じデータでこれらを選ぶと，数値が水増しされる．
% そこで研究を2段階に分けた：見ることを許した探索段階と，見なかった最終テストである．
% 探索段階は 8 層 × 3 種の方向を 100 プロンプトで比べ，最良の組（第4層・プローブ方向）を同じプロンプト——プローブの訓練にも使われた——で採点した；この探索は補正していないので，この段階は探索として報告し，テストとしては報告しない．
% 最終テストは，プローブもそれ以前のどの実行も見ていないテスト分割の領域から新規100プロンプトを引き，同一調の標的を除くと1条件 1,100 本になる．
% 層や方向からシードや標本サイズまでの全選択は，この段階の最初の数値が出る前に固定し，計画からのずれはすべて報告する（計画そのものは補足資料）．
% 副条件は安定した短調プロンプトと短調12標的で同じ一式を繰り返す；その設計も短調の実行が1つも無い時点で固定したので，通るか落ちるかしかない．
The procedure of Section~\ref{sec:exp2} leaves choices open---layer, directions, positions---and settling them on the data that produces the reported numbers inflates them.
We therefore ran two stages: a \emph{search stage} allowed to look ($8$ layers $\times$ $3$ direction kinds on $100$ prompts the probe had also trained on; uncorrected, so reported as a search, not a test), and a \emph{final test} of $100$ fresh prompts from a region of the test split no earlier run had seen---for each prompt the $11$ targets other than its own key, so $100\times11 = 1{,}100$ continuations per condition.
Every choice, from the layer down to seeds and sample size, was fixed before this stage's first number existed; departures are reported, and the supplementary gives the plan.
A secondary condition repeats the battery with stable-minor prompts and minor targets, its design fixed before any minor run existed, so it could only pass or fail.

% IV. 結果
\section{Results}
\label{sec:results}
% 2つの問いを順に扱う：まず調がそこにあるか（4-1 節），次にモデルがそれを使っているか（4-2–4-4 節）．最後に，どちらの結果も実音楽で生き残るかを確かめる（4-5 節）．
% 報告する数値は移調増強あり・シード0のモデルのものである．
We ask first whether the key is there (Sec.~\ref{sec:h1}), then whether the model uses it (Secs.~\ref{sec:h3}--\ref{sec:h5}), then whether either survives contact with real music (Sec.~\ref{sec:wild}); numbers are from the augmented seed-$0$ model.

\subsection{Experiment 1: The Key Is Easier to Read}
\label{sec:h1}
% モデルは調を自ら計算しており，入力の表面から写しているのではない．
% 第4層でプローブは macro-F1 0.927 に達し，最良の音符数えベースラインは 0.790 にとどまる：control task の床 0.032 0.032 を引くと，プローブの優位は +0.105 である（ブートストラップ 95% CI [0.091, 0.114]）．
% この優位は特定のネットワークや特定の深さの癖ではない：訓練済み6モデルすべての第1–7層で，信頼区間はゼロを上回ったままである．
% 例外は第0層だけである（−0.20）：入力の直後なので，活性は音符をほぼそのまま写しており，音符数えのほうがよく読む．未訓練のネットワークではプローブは 0.243 にとどまる．
% 窓を広げても音符数えは追いつかない：最強の拡張版が 0.824 で，プローブの優位は +0.071 [0.056, 0.081] 残る．
% 我々の主張は，これらの音符数えより線形に読みやすいということまでである．音符から調を計算するあらゆる方法に勝つとは主張しない．
The model computes the key; it does not copy it from the surface of the input.
At layer~$4$ the probe reaches macro-$F_1$ $0.927$ against $0.790$ for the best note-counting baseline; subtracting the control-task floor of $0.032$ leaves a margin of $+0.105$ ($95\%$ CI $[0.091,0.114]$), and the CI stays above zero in layers $1$--$7$ of all six trained models.
The one exception, layer~$0$ ($-0.20$), sits just after the input, where the activation still mirrors the notes; an untrained network reaches only $0.243$.
The wider windows do not help the note counter catch up (margin $+0.071$ $[0.056,0.081]$ against the strongest); we claim only that the key is easier to read linearly from the activation than from these note counts.
% [SUPPL] 拡張ベースラインの内訳と窓ごとのスコア。窓が転調をまたぐと，転調前の音符が数え上げを古い調へ引き戻す点も補足で述べる。

\subsection{Experiment 2: The Music Follows the Installed Key}
\label{sec:h3}
% 音楽は，我々が書き込んだ調へ移る．
% 最終テストでは，およそ3本に1本の継続が書き込んだ調にちょうど着地する（SR 0.355）のに対し，同じ編集をランダム方向へ行うと着地はおよそ 25 本に1本（0.039）である——比 9.1（表1）．
% 効果は特定の調に限られない：12の標的調すべてで編集がベースラインを上回り，12調にわたる Holm 補正の後もそれぞれ個別に有意である（最悪 p=4.8×10^-5，r=0.72–1.00）；プロンプト単位の差のブートストラップ 95% CI は [0.284, 0.347] である．
The music moves into the key we install: in the final test roughly one continuation in three lands in exactly the installed key ($0.355$) against one in twenty-five for random directions ($0.039$), a ratio of $9.1$ (Table~\ref{tab:main}).
The edit beats the baseline for every target key, each significant after Holm correction (worst $p=4.8\times10^{-5}$, $r=0.72$--$1.00$); the prompt-paired difference has a $95\%$ CI of $[0.284,0.347]$.

% 表1：最終テストでの編集の成功率．
% 上段が主結果，中段は対抗説明を消す比較，下段は上限と，何も変えない場合の値である．
\begin{table}[!t]
\centering
\footnotesize
\setlength{\tabcolsep}{4pt}
\caption{Success rate of the edit on the final test.
Top: the main result. Middle: comparisons that rule out other explanations. Bottom: the ceiling, and what happens when we change nothing.}
\label{tab:main}
\begin{tabularx}{\columnwidth}{@{}Xcc@{}}
\toprule
 & \multicolumn{2}{c}{SR}\\
\cmidrule(l){2-3}
What we write & major & minor\\
\midrule
\textbf{The key we install}, at every position & \textbf{0.355} & \textbf{0.495}\\
\midrule
\emph{Directions}: random, rank-matched & 0.039 & 0.003\\
\emph{Distance}: random, rank- and norm-matched & 0.056 & 0.026\\
\emph{Positions}: the key at pitch positions only & 0.035 & 0.000\\
\quad at bar and note-length positions only & 0.233 & 0.212\\
\emph{Operation}: the key added, norm-matched & 0.245 & 0.154\\
\quad added at twice the displacement & 0.471 & 0.505\\
\quad added as a difference of means & 0.386 & 0.435\\
\midrule
Transposing the prompt instead of editing & 0.648 & 0.877\\
The key the prompt is already in & 0.650 & 0.800\\
\bottomrule
\end{tabularx}
\end{table}

% 推定器に依存しているのでは？ 違う：推定器を使わずに音符を数えても同じ結論になる．
% 継続の音のうち，書き込んだ調の音階に属するものは 0.936，プロンプトの調の音階に属するものは 0.617 である．
\noindent\textbf{An artifact of the key estimator?}
No: with no estimator at all, $0.936$ of continuation notes belong to the installed key's scale, $0.617$ to the prompt key's.

% 「音楽が壊れていない」という条件が差を作ったのでは？ 違う：条件 (ii) を外すと，編集は 0.355 から 0.410 に上がるのに対し，ランダムベースラインは 0.039 と 0.056 から 0.041 と 0.059 にしか上がらない．
% 削られていたのは編集のほうであり，この条件は差を作るどころか縮めていた．
\noindent\textbf{A gift from the ``not broken'' condition?}
No: dropping condition~(ii) raises the edit from $0.355$ to $0.410$ and the random baselines only to $0.041$ and $0.059$ --- the condition narrowed the gap, not created it.

% より強い押しでは？ 違う：編集と同じ本数の方向を，同じ距離だけ動かすランダムベースラインでも，0.056 にしか届かない；編集は 12 の調すべてでこれを上回る．
% 重要なのはどの方向に書くかであって，どれだけ強く押すかではない．
\noindent\textbf{A harder push?}
No: matching the random write to the edit's distance as well still reaches only $0.056$, beaten for all $12$ keys --- what matters is \emph{which} directions we write, not how far.

% 使われていない経路による錯覚では？ 部分空間へのパッチは，モデルが通常使わない経路を通して出力を動かすことがある．
% ここでは編集はモデル自身の符号化に従って働く：音符が生成される前に確率を動かし（IV-C 節），小節と音長の位置でだけ効く（IV-D 節）．
\noindent\textbf{An illusion from a dormant pathway~\cite{makelov2024illusion}?}
The edit follows the model's own encoding instead: it moves probabilities before any note exists (Sec.~\ref{sec:nextpitch}) and acts only at bar and note-length positions (Sec.~\ref{sec:h5}).

\noindent\textbf{How large is $0.355$?}
% プロンプト自体の移調——出力が起こしうる最大の移動——は 0.648 に達する；これを移調上限と呼び，編集はその 55% を覆う．決して触らないプロンプトの音符が継続を古い調へ引き戻し続けるので，編集はこの上限より下に留まる．
% 調を書き込むこと自体が音楽を壊すわけではない：プロンプトがすでにいる調を書き込むと，モデルはその調に留まり（0.650，移調上限と同水準），ランダムな方向へ書き込んだ場合は 51.7% しか留まらない．
Transposing the prompt itself, the largest move the output can make, reaches $0.648$; we call this the \emph{transposition ceiling}, and the edit covers $55\%$ of it, held below by the prompt's own notes, which we never touch and which keep pulling the continuation back toward the old key.
Writing a key is not in itself disruptive: install the key the prompt is already in and the model stays in it ($0.650$, level with the ceiling).
% [SUPPL] 条件 (ii) を適用する前の identity 値 0.710。何もしない編集（健全性検査）は 100/100 プロンプトで編集なし出力をトークン単位で再現。

% 短調でも同じことが起こる：事前に固定した副条件（III-D 節）で，編集は 0.003 に対して 0.495 に達し，12 の標的調すべてで有意である（最悪 p=1.7×10^-7，r≥0.96）．
\noindent\textbf{Minor keys.}
The pre-registered minor condition (Sec.~\ref{sec:confirm}) repeats everything, larger (Table~\ref{tab:main}, right column; all $12$ keys significant, worst $p=1.7\times10^{-7}$, $r\ge0.96$).

\subsection{The Edit Acts First, the Notes Carry It}
\label{sec:nextpitch}
% 編集は，モデルが音符を1つも生成しないうちに次の音の確率を動かす．
% 書き込んだ調の音階上の確率質量とプロンプトの調の音階上の質量との対数比を 0.695 上げる——およそ2倍である——のに対し，ランダムベースラインは 0.041 しか上げない．
% 差は補正後も 12 の調すべてで有意である（最悪 p=5×10^-16，r≥0.99）．
% したがって，モデルが新しい調を自分の出力から読んでいるという説明は成り立たない．編集の効果は，読み取る出力が存在する前に生じている．
The edit moves the next-note probabilities before the model generates a single note: the log ratio of probability mass on the installed key's scale over the prompt key's rises by $0.695$---roughly doubling---against $0.041$ for the random baseline, significant for all $12$ keys (worst $p=5\times10^{-16}$, $r\ge0.99$).
The model therefore cannot be reading the new key from its own output: the effect precedes any output.

% 最初の1手の後は，音符が調を運ぶ．モデルが我々の書いた値を覚えているわけではない．
% 探索段階で，一度だけ書き込んで1小節を生成させ，その1小節のトークンを同じプロンプトに継ぎ足し，編集を一切効かせずに続きを書かせると，全位置に書き込んだ場合と同じ効果が出た．
% モデルの内部には何も残っていないのに効果は残る．
After that first decision the notes carry the key, not a memory: writing once, letting the model generate one bar, then appending that bar's tokens to the same prompt and continuing with no edit anywhere matched writing at every position (search stage).
Nothing remained inside the model, yet the effect did.

\subsection{Where the Key Is Read and Where It Acts}
\label{sec:h5}
% 調は，使えるようになる前に読めるようになる．
% 深さ方向では，探索段階で，生のプローブ得点は第2–6層で平ら（macro-F1 0.920–0.929）だが，編集が効くのは第2層以降だけである：第2層以降ではランダムベースラインを 0.19–0.30 上回るのに対し，プローブがすでに調をよく読む第0層と第1層では 0.03 と 0.04 にとどまる．
The key becomes readable before it becomes usable: the raw probe score is flat across layers~$2$--$6$ ($0.920$--$0.929$), yet the edit beats its baseline by $0.19$--$0.30$ only from layer~$2$ on, against $0.03$--$0.04$ at layers~$0$--$1$, where the probe already reads the key well (search stage).

% 編集が効くのは，モデルが音高を出す位置ではなく，小節と音長を担う位置である．
% 音高位置だけへの書き込みは何もしない：0.035 で，ランダムベースラインの 0.039 を下回り，有意になる調は1つもない．
% 小節と音長の位置だけへの書き込みは，全効果の 65% にあたる 0.233 に達し，12 の調のうち 11 で有意である．
% 短調条件でも同じことが起きる：音高位置のみは 0.000，小節と音長の位置のみは全効果の 43% にあたる 0.212 である．
% 音高位置のみの編集は，全位置への編集（84%）より高い割合で条件 (ii) を満たす（98%）が，それは音楽をほとんど変えないからにすぎない．
% したがって，編集が出力を変えるのは位置の一部だけである．
The edit acts at the positions that carry bars and note lengths, not where the model emits a pitch (Table~\ref{tab:main}, middle).
Writing only at pitch positions does nothing ($0.035$, below the random baseline, no key significant); writing only at bar and note-length positions carries $65\%$ of the full effect ($11/12$ keys), and the minor condition repeats both.
Pitch-only edits pass condition~(ii) more often ($98\%$ against $84\%$) only because they barely change the music.
% [SUPPL] 層をまたぐ書き込みの観察: 第5–7層の方向を4層手前に書くと 0.27–0.39 で働き，逆向きは 0.003–0.082（探索段階のランダムベースラインは 0.075）。方向は深さをまたいで通用するが，出力を変えるのは狭い帯の深さだけ。

\subsection{A Public Model}
\label{sec:wild}
% 主セル（AMT×バッハ）：プローブ優位 +0.196、編集はランダムの 7.9 倍．サイズ軸（2倍・6倍）では何も変わらない．
% ジャンル軸とトークン化軸では読み取りと使用が「二重に」解離する：
%  AMT×pop：読める（+0.132）が編集は事前登録基準で不成立（0.135 対 0.057、1/12）．
%  MMT×pop：プローブは音符数えに負ける（−0.061）のに、同じ部分空間の編集が 12/12 で成立（0.360 対 0.064、IKR 交差 0.961/0.660）．
% すなわち可読性は因果的使用の必要条件でも十分条件でもない．
% 十字形の設計を本文で先に一文で述べる（表を見る前に読者が空欄の意味を理解できるように）．
The grid is a cross, not a product: we fix the model and change the corpus, then fix the corpus and change the tokenization, because a Bach cell for the pop-trained checkpoints would confound tokenization with distribution shift.
In the primary cell (Table~\ref{tab:wild}) the probe beats the note-counting baselines by $+0.196$ $[0.132,0.260]$ and the edit reaches $7.9\times$ the random baseline, with $99.7\%$ of continuations meeting condition~(ii) and the notes crossing over ($0.890$ against $0.753$); doubling or sextupling the model changes none of it (margins $+0.196/+0.213/+0.203$, all CIs overlapping; the $24$-layer edit at $0.557$), and the acting layer sits above the best-reading layer in every case.
Crossing genre and tokenization, reading and using come apart in BOTH directions.
The Anticipatory model still reads the key on pop ($+0.132$) yet its edit fails the pre-registered bar ($0.135$ against $0.057$; one key of twelve).
MMT's probe LOSES to the note counts on the same corpus ($-0.061$ $[-0.13,-0.003]$), yet editing the same probe-derived subspace moves its music decisively: $0.360$ against $0.064$, all $12$ keys, with the in-key shares crossing ($0.961$ against $0.660$).
% 可読性は因果的使用の必要条件でも十分条件でもない．ここから出るのは「プローブでは決着しない」（十分性の主張）であり，
% 「上書きだけが唯一の手段」ではない．上書きの固有の強みは特定の値を「設定」できることにある．
Readability is therefore neither necessary nor sufficient for causal use: a probe cannot settle the question either way, while the overwrite settles it --- and, unlike ablating or nudging, it is the test that can \emph{set} a specific value.

% 表2：公開モデル×コーパス．プローブ優位（補正後、95% CI）と編集の成功率（対ランダム）．
\begin{table}[!t]
\centering
\footnotesize
\setlength{\tabcolsep}{3.5pt}
\caption{Public models. Probe margin: control-corrected macro-$F_1$ over the best note-counting baseline. Edit: guarded success rate against the rank-matched random write; every cell without $\dagger$ separates for all $12$ target keys after Holm correction. $\dagger$ marks pre-registered FAILURES: the AMT pop edit reaches only $1/12$ keys, and the MMT pop probe loses to the note counts (CI $[-0.13,-0.003]$). All checkpoints were trained by their authors on Lakh MIDI; none ever saw our data. A dash is a cell we did not run, and deliberately: the design varies ONE thing at a time --- model fixed while the corpus changes, corpus fixed while the tokenization changes --- and MMT and the REMI baseline were trained on pop-heavy web MIDI, so a Bach cell for them would confound tokenization with distribution shift.}
\label{tab:wild}
\begin{tabularx}{\columnwidth}{@{}Xllll@{}}
\toprule
 & \multicolumn{2}{l}{Bach chorales} & \multicolumn{2}{l}{Pop (POP909-CL)}\\
\cmidrule(lr){2-3}\cmidrule(l){4-5}
Model & probe & edit & probe & edit\\
\midrule
AMT small (12L) & $+.196$ & $.479$ vs $.061$ & $+.132$ & $.135$ vs $.057^{\dagger}$\\
AMT medium (24L) & $+.213$ & $.557$ vs $.057$ & --- & ---\\
AMT large (36L) & $+.203$ & --- & --- & ---\\
MMT (compound) & --- & --- & $-.061^{\dagger}$ & $.360$ vs $.064$\\
REMI repr. & --- & --- & TBD & TBD\\
\bottomrule
\end{tabularx}
\end{table}

% V. 考察
\section{Discussion}
\label{sec:discussion}
% 我々が測った分かれ方は，方向をどう選ぶかに関わる．
% 言語モデルも同じ隔たりを報告している．
% したがって方向は，読むよりも書くことで検証したほうがよい．
The reading/acting split (Sec.~\ref{sec:h5}) bears on how steering directions are chosen: music steering~\cite{smitin2024,facchiano2025patching,singh2026discovering,musicrfm2026} selects them by how well an attribute decodes, and our pitch-position condition shows that criterion picking a site the model does not use --- the gap language models also report~\cite{wu2025axbench,arad2025selecting}.
A direction is therefore better validated by writing to it than by reading from it, and we expect the same gap for other attributes recovered from music activations.
% 追加：同一パイプライン上での install 対 steering の直接比較（事前登録済み・実行中）。
% 加算は部分空間内の旧調成分 P_V h を保持するので（P_V h' = P_V h + δu が解析的に成り立つ）、
% 置換との差は「旧成分を消すか残すか」だけに縮まる．結果は表1の追加行（TBD）に入る．
We also pre-registered a direct comparison on this pipeline: the same value \emph{added} (steering) instead of installed by Eq.~\eqref{eq:edit}, matched per position to the install edit's own displacement norm, so the only remaining difference is whether the old key component $P_Vh$ is removed --- addition provably preserves it.
% held-out の確定結果（表1下段）：
%  同変位量の加算 0.245 対 置換 0.355（差 +0.111、CI [+0.078,+0.145]、12標的中3で個別有意）→ 事前登録の向きで支持．
%  2倍変位の加算 0.471 → 事前登録と逆向き（CI [-0.154,-0.078]）．そのまま報告する．
%  内訳が機序を示す：同変位量の加算では継続の 28.8% がプロンプトの調に留まる（置換 6%台）．差分平均方向では 0.6%．
%  つまり「加算では設定できない」ではなく「摂動あたりの効率で置換が優り，旧成分の除去が効いている」．
We ran that comparison on the held-out prompts, paired by prompt against the install rows (Table~\ref{tab:main}, bottom).
At matched displacement the replacement wins, by $0.111$ ($95\%$ CI $[0.078,0.145]$; three of twelve keys individually significant after correction), and the estimator says why: $28.8\%$ of the addition's continuations stay in the PROMPT's key, which is the preserved $P_Vh$ showing up in the notes.
Allowed twice that displacement, however, the addition overtakes the replacement ($0.471$ against $0.355$) --- the opposite of our pre-registered direction, reported as it fell --- and it pays for the extra push at $4\times$, where half the continuations break the quality budget, and at $8\times$, where nearly all do.
The difference-of-means direction leaves the prompt key almost nowhere ($0.6\%$ of continuations, in-key share $0.438$), which is the removal of the old component happening implicitly rather than by subtraction.
So addition can set a key; what replacement buys is efficiency per unit of perturbation, and what the whole comparison isolates is that removing what is already written is doing part of the work.

\section{Limitations}
\label{sec:limits}
% 我々が示したのは，モデルが調を計算し，生成しながら使っていることである．人間のように調を理解しているとは主張しない．
% 条件 (ii) が測るのは，編集が基準モデルをどれだけ乱すかであって，音楽としての良さではない．聴取実験は行っていない．
% 公開モデルで確かめたのはモデル1つだけである．
We show that the model computes a key and uses it as it generates; we do not claim that it understands keys the way a person does.
The dormant-pathway controls of Secs.~\ref{sec:nextpitch} and~\ref{sec:h5} were run on our own models, not on the public ones, so the MMT edit result is a behavioral finding whose internal pathway we have not localized.
Condition~(ii) measures how far an edit disturbs a reference model, not musical quality itself, and we ran no listening test.
% 公開モデル検証は「1モデル」から「3サイズ＋2トークン化＋2ジャンル」へ拡張中（TBD セルは実行中）．
The public-model check now spans three checkpoint sizes of one family, two further tokenizations, and a second genre; the cells marked TBD in Table~\ref{tab:wild} were still running when this draft was built, and any that fail will be reported as failing.
% [SUPPL] 公開モデルでは層を探索段階で選んだままで，新規プロンプトでの確認は行っていない。

\section{Conclusion}
% 我々が測れた範囲では，音楽 Transformer は自分がどの調にいるかを知っている：モデルは調を計算し，我々が置き換えられる形で保持し，生成しながらそれを参照している．
% 読めることは，使っていることの証拠にはならない．
% テストベッドとコード，そしてデモを公開する．
In the sense our test can reach, a music Transformer does know what key it is in: it computes a key, holds it in a form we can replace, and consults it as it generates---while being readable, alone, is not evidence of being used.
We release the testbed, code, and demos: \url{https://github.com/<user>/<repo>}.

\bibliographystyle{IEEEbib}
\bibliography{refs}

\end{document}
